\newpage
\pagenumbering{roman}
\setcounter{page}{1}
\section{Glossar der Abkürzungen}

Dieses Glossar wurde mit der Hilfe von Künstlicher Intelligenz ausgearbeitet.

\begin{itemize}
    \item \textbf{1ACD} -- 1st Aircat Division: Bezeichnung der Arma-3-Unit, auf die sich das Regelwerk bezieht.
    \item \textbf{a.D.} -- außer Dienst: Zusatz zur Kennzeichnung von Mitgliedern, die sich im Status „außer Dienst“ befinden.
    \item \textbf{d.R.} -- der Reserve: Zusatz zur Kennzeichnung von Mitgliedern, die sich in der Reserve befinden.
    \item \textbf{Hptm.} -- Hauptmann: Offiziersdienstgrad, der innerhalb der Unit verwendet wird.
    \item \textbf{OLt.} -- Oberleutnant: Offiziersdienstgrad, der innerhalb der Unit verwendet wird.
    \item \textbf{Lt.} -- Leutnant: Offiziersdienstgrad, der innerhalb der Unit verwendet wird.
    \item \textbf{OStFw.} -- Oberstabsfeldwebel: Unteroffiziersdienstgrad mit Portepee, der innerhalb der Unit verwendet wird.
    \item \textbf{StFw.} -- Stabsfeldwebel: Unteroffiziersdienstgrad mit Portepee, der innerhalb der Unit verwendet wird.
    \item \textbf{HptFw.} -- Hauptfeldwebel: Unteroffiziersdienstgrad mit Portepee, der innerhalb der Unit verwendet wird.
    \item \textbf{OFw.} -- Oberfeldwebel: Unteroffiziersdienstgrad mit Portepee, der innerhalb der Unit verwendet wird.
    \item \textbf{Fw.} -- Feldwebel: Unteroffiziersdienstgrad mit Portepee, der innerhalb der Unit verwendet wird.
    \item \textbf{StUffz.} -- Stabsunteroffizier: Unteroffiziersdienstgrad ohne Portepee, der innerhalb der Unit verwendet wird.
    \item \textbf{Uffz.} -- Unteroffizier: Unteroffiziersdienstgrad ohne Portepee, der innerhalb der Unit verwendet wird.
    \item \textbf{StKorp.} -- Stabskorporal: Mannschaftsdienstgrad, der innerhalb der Unit verwendet wird.
    \item \textbf{Korp.} -- Korporal: Mannschaftsdienstgrad, der innerhalb der Unit verwendet wird.
    \item \textbf{OStGefr.} -- Oberstabsgefreiter: Mannschaftsdienstgrad, der innerhalb der Unit verwendet wird.
    \item \textbf{StGefr.} -- Stabsgefreiter: Mannschaftsdienstgrad, der innerhalb der Unit verwendet wird.
    \item \textbf{HptGefr.} -- Hauptgefreiter: Mannschaftsdienstgrad, der innerhalb der Unit verwendet wird.
    \item \textbf{OGefr.} -- Obergefreiter: Mannschaftsdienstgrad, der innerhalb der Unit verwendet wird.
    \item \textbf{Gefr.} -- Gefreiter: Mannschaftsdienstgrad, der innerhalb der Unit verwendet wird.
    \item \textbf{Jg.} -- Jäger: Mitglied auf Probe innerhalb der Unit.
    \item \textbf{PTS} -- Permission to Speak: Begriff, der während einer Mission oder Besprechung genutzt wird, um das Recht zu erhalten, sprechen zu dürfen.
    \item \textbf{Unit-Leiter} -- Der Unit-Leiter ist die führende Person innerhalb der 1st Aircat Division. Diese Person hat die oberste Verantwortung für alle administrativen und organisatorischen Aufgaben. Der Unit-Leiter trifft wichtige strategische Entscheidungen und leitet die Unit in allen Angelegenheiten. Der Unit-Leiter wird durch die Mitgliederversammlung gewählt.
    \item \textbf{Spieß} -- Der Spieß ist der Stellvertreter des Unit-Leiters und hat insbesondere eine unterstützende Rolle in administrativen und organisatorischen Aufgaben. Der Spieß übernimmt die Aufgaben des Unit-Leiters im Falle von Abwesenheit des Unit-Leiters. Der Spieß wird ebenfalls durch die Mitgliederversammlung gewählt.
    \item \textbf{GvD} -- Gefreiter vom Dienst: Ein Mannschaftsvertreter, der während seiner Amtszeit die Interessen der Mannschaft gegenüber dem Stab und erweiterten Stab vertritt.
    \item \textbf{UvD} -- Unteroffizier vom Dienst: Ein Unteroffizier, der während seiner Amtszeit bestimmte administrative Aufgaben übernimmt und die Interessen der Unteroffiziere vertritt. Mitglied des Stabs.
    \item \textbf{S1-Feldwebel} -- Feldwebel der inneren Führung: Zuständig für die Kontrolle und Koordination der Leitung sowie oberste Beschwerdeinstanz. Er übernimmt die Leitung der Unit in Abwesenheit der Leitung. Mitglied des erweiterten Stabs.
    \item \textbf{PersFw} -- Personal-Feldwebel: Zuständig für die Unterstützung des Gefreiten vom Dienst und die Kontrolle der Mitgliederverwaltung. Mitglied des erweiterten Stabs.
    \item \textbf{S6-Feldwebel} -- Feldwebel der Technik: Zuständig für technische Fragen und die Verwaltung öffentlicher Kanäle der Unit. Mitglied des erweiterten Stabs.
    \item \textbf{Leitung} -- Die Leitung setzt sich aus dem Unit-Leiter und dem Spieß zusammen. Sie trägt die Verantwortung für die strategische Ausrichtung und die täglichen administrativen Aufgaben der Unit. Entscheidungen der Leitung gelten sofort, bis sie durch die Mitgliederversammlung oder eine neue Entscheidung wieder aufgehoben werden.
    \item \textbf{Stab} -- Der Stab unterstützt die Leitung der Unit in administrativen und organisatorischen Aufgaben. Der Stab besteht aus der Leitung (Unit-Leiter und Spieß) sowie weiteren Positionen wie dem Gefreiten vom Dienst (GvD) und dem Unteroffizier vom Dienst (UvD). Der Stab tagt regelmäßig, mindestens einmal im Quartal.
    \item \textbf{erweiterter Stab} -- Der erweiterte Stab umfasst alle Mitglieder des Stabs sowie alle aktiven Unteroffiziere mit und ohne Portepee. Der erweiterte Stab unterstützt die Leitung bei strategischen und administrativen Entscheidungen und tritt in regelmäßigen Sitzungen zusammen. Entscheidungen im erweiterten Stab betreffen in der Regel wichtige organisatorische Maßnahmen und die langfristige Ausrichtung der Unit sowie Streitbeilegungen.
    \item \textbf{Zeus} -- Missionsbauer/Spielleiter in Arma 3: Rolle in Arma 3 zur Leitung oder Erstellung von Missionen.
\end{itemize}
